\documentclass{article}
\usepackage[utf8]{inputenc}
\usepackage[T1]{fontenc}
\usepackage{geometry}
\geometry{
    a4paper,
    left=2.5cm,
    right=2.5cm,
    top=2.5cm,
    bottom=2.5cm
}
\usepackage{graphicx}
\usepackage{float}
\usepackage{hyperref}
\usepackage[english]{babel}
\usepackage{amsmath}
\usepackage{amssymb}
\usepackage{array}
\usepackage{booktabs}
\usepackage{caption}
\usepackage{subcaption}
\usepackage{enumitem}
\usepackage{fancyhdr}
\usepackage{titlesec}
\usepackage{url}
\usepackage{textcomp}
\usepackage{xcolor}
\usepackage{listings}
\usepackage{multirow}
\usepackage{longtable}
\usepackage{tabularx}
\usepackage{times}

% Set up page style
\pagestyle{fancy}
\fancyhf{}
\rfoot{\thepage}
\renewcommand{\headrulewidth}{0pt}

% Title formatting
\titleformat{\chapter}[display]{\normalfont\bfseries}{\chaptertitle\thechapter}{0pt}{\Huge}
\titlespacing*{\chapter}{0pt}{50pt}{40pt}

% Document info
\title{Template Document}
\author{Academic Support Agent}
\date{\today}

\begin{document}
\maketitle

\pagestyle{meuestilo}
\checkandfixthelayout
% Seleciona o idioma do documento (conforme pacotes do babel)
%\selectlanguage{english}
\selectlanguage{brazil}

% Retira espa o extra obsoleto entre as frases.
\frenchspacing 

% ----------------------------------------------------------
% ELEMENTOS PR -TEXTUAIS
% ----------------------------------------------------------

%---
%
% Se desejar escrever o artigo em duas colunas, descomente a linha abaixo
% e a linha com o texto ``FIM DE ARTIGO EM DUAS COLUNAS''.
% \twocolumn[       % INICIO DE ARTIGO EM DUAS COLUNAS
%
%---

% p gina de titulo principal (obrigat rio)
%\maketitle




% titulo em outro idioma (opcional)



% ----------------------------------------------------------
% ELEMENTOS TEXTUAIS
% ----------------------------------------------------------
\textual
\pagestyle{meuestilo}

\phantom{This text will be invisible}
\vfill

\begin{center}
\imprimirtitulo
\end{center}

\vfill

\begin{center}
Local de execu o: \imprimirlocal
\end{center}

\vfill

\begin{center}
\imprimirdata
\end{center}
\newpage

\begin{center}
\begin{table}[h]
%\begin{tabular}{p{.3cm}p{1.9cm}p{1.1cm}p{3.9cm}p{0.3cm}p{1.1cm}p{5.1cm}}
\begin{tabularx}{\linewidth}{aeioaeu}

 &  & & & & &  \\\hline
\multicolumn{7}{|l|}{\textbf{1 -- DADOS DO PROJETO}}                                                                                             \\ \hline
\multicolumn{7}{|l|}{T tulo: \imprimirtitulo}                                \\ \hline
\multicolumn{7}{|l|}{ rea Tem tica Principal}                                                                                                   \\ \hline
%\multicolumn{7}{|l|}{\textit{(assinalar apenas uma  rea   Tem tica, aquela que melhor define o problema que deu origem ao projeto)}} \\ \hline
\multicolumn{1}{|c|}{}        & \multicolumn{3}{l|}{Comunica o}    & \multicolumn{1}{c|}{}          & \multicolumn{2}{l|}{Meio Ambiente}     \\ \hline
\multicolumn{1}{|c|}{}         & \multicolumn{3}{l|}{Cultura}        & \multicolumn{1}{c|}{}           & \multicolumn{2}{l|}{Sa de}             \\ \hline
\multicolumn{1}{|c|}{} &
	\multicolumn{3}{l|}{Direitos Humanos e Justi a} &
	\multicolumn{1}{c|}{X} &
	\multicolumn{2}{l|}{Tecnologia e Produ o} \\ \hline
\multicolumn{1}{|c|}{}         & \multicolumn{3}{l|}{Educa o}       & \multicolumn{1}{c|}{}           & \multicolumn{2}{l|}{Trabalho}          \\ \hline
\multicolumn{3}{|l|}{\begin{tabular}[c]{@{}l@{}}Local de realiza o/\\ Munic pio\end{tabular}} &
	\multicolumn{4}{l|}{Microregi o de SUAPE no Estado de Pernambuco, Brasil} \\ \hline
\multicolumn{3}{|l|}{Comunidade}                       & \multicolumn{4}{l|}{Comunidades dos munic pios de Ipojuca e Cabo de Santo Agostinho} \\ \hline
\multicolumn{2}{|l|}{\begin{tabular}[c]{@{}l@{}}Benefici rios \\ Diretos\end{tabular}} &
	\multicolumn{2}{l|}{\begin{tabular}[c]{@{}l@{}}Moradores da \\ regi o de X\end{tabular}} &
	\multicolumn{2}{l|}{\begin{tabular}[c]{@{}l@{}}Benefici rios \\ Indiretos\end{tabular}} &
	\multicolumn{1}{l|}{\begin{tabular}[c]{@{}l@{}}Trabalhadores da \\ empresa X\end{tabular}} \\ \hline
\end{tabularx}
\end{table}
\end{center}

\begin{table}[h]
		\centering
		\begin{tabularx}{\textwidth}{|g|g|s|s|s|}
\hline
\multicolumn{5}{|l|}{\textbf{2 -- INSTITUI ES PARCEIRAS}}                                                                                             \\ \hline


{\multirow{2}{*}{Nome completo da Institui o}} &
	{\multirow{2}{*}{Fun o}} &
	\multicolumn{3}{c|}{Tipo} \\ \cline{3-5} 
\multicolumn{1}{|c|}{} &
	\multicolumn{1}{c|}{} &
	\multicolumn{1}{c|}{P blica} &
	\multicolumn{1}{c|}{Privada} &
	\multicolumn{1}{c|}{ONG} \\ \hline
Casa da M e Joana &
	Fornecer abrigo &
	\multicolumn{1}{c|}{\textbf{X}} &
	\multicolumn{1}{c|}{\textbf{}} &
	\multicolumn{1}{c|}{\textbf{}} \\ \hline
F brica de Livros &
	Imprimir livros &
	\multicolumn{1}{c|}{\textbf{}} &
	\multicolumn{1}{c|}{\textbf{X}} &
	\multicolumn{1}{c|}{\textbf{}} \\ \hline
Amigos do ambiente &
	vas &
	\multicolumn{1}{l|}{\textbf{}} &
	\multicolumn{1}{l|}{\textbf{}} &
	\textbf{} \\ \hline
		\end{tabularx}
\end{table}

\begin{theo}[APRESENTA O]
    Este   um modelo para o template de submiss o de PIBEX do IFPE \cite{abntex2-wiki-como-customizar}. Aqui tem um exemplo de refer ncia bibliogr fica \cite{ibge1993}.
    
    \lipsum[1-3]
\end{theo}

\begin{theo}[CARACTERIZA O E JUSTIFICATIVA]
    \lipsum[4-6]
\end{theo}

\begin{theo}[OBJETIVOS]
	\lipsum[7]
\end{theo}

\begin{theo}[METODOLOGIA]
	\lipsum[8-12]
\end{theo}

\begin{theo}[RESULTADOS ESPERADOS]
	\lipsum[13-14]
\end{theo}

\begin{theo}[ACOMPANHAMENTO E AVALIA O]
	\lipsum[15]
\end{theo}

\begin{theo}[CRONOGRAMA]
	\lipsum[16]
\end{theo}

\begin{theo}[RECURSOS E INFRAESTRUTURA EXISTENTES]
	\lipsum[17-18]
\end{theo}

\begin{theo}[RECURSOS E INFRAESTRUTURA NECESS RIOS]
	\lipsum[19-20]
\end{theo}

\renewcommand{\bibname}{}

\begin{theo}[REFER NCIAS]
\vspace{-1.6cm} 

\end{theo}
\end{document}